\chapter{PAC Learning}

\section{Scope Of This Chapter}

This chapter fixes the following scope:

\begin{itemize}
\item binary classification in the main body of the chapter
\item distribution-free PAC learning
\item polynomial sample complexity as the baseline statistical notion
\item polynomial running time as the computational refinement
\item proper and improper learners treated as distinct notions
\end{itemize}

The guiding question is:

\begin{quote}
If realizable PAC learning and agnostic PAC learning are statistically equivalent, when does that equivalence survive once we also demand polynomial-time algorithms?
\end{quote}

A later section records how this picture changes in multiclass classification and in regression with continuum labels.

\section{Definitions}

Let $\mathcal{X}$ be an instance space and let $\mathcal{C} \subseteq \{0,1\}^{\mathcal{X}}$ be a concept class.

\paragraph{Realizable PAC learning.}
A learner receives i.i.d.\ samples $(x,c(x))$ where $x \sim D$ for an arbitrary distribution $D$ over $\mathcal{X}$ and an unknown target concept $c \in \mathcal{C}$. It must output $h$ such that
\[
\Pr\!\left[\err_D(h,c) \le \varepsilon\right] \ge 1-\delta.
\]

\paragraph{Agnostic PAC learning.}
A learner receives i.i.d.\ samples from an arbitrary distribution $\mathcal{D}$ over $\mathcal{X} \times \{0,1\}$ and must output $h$ such that
\[
\Pr\!\left[\err_{\mathcal{D}}(h) \le \inf_{h' \in \mathcal{H}} \err_{\mathcal{D}}(h') + \varepsilon\right] \ge 1-\delta.
\]
Here $\mathcal{H}$ is the benchmark hypothesis class.

\paragraph{Proper vs improper.}
The learner is \emph{proper} if it must output $h \in \mathcal{C}$. It is \emph{improper} if it may output $h$ from a larger class $\mathcal{H}$.

\paragraph{Efficient learning.}
In this document, ``efficient'' means both:
\begin{itemize}
\item sample complexity polynomial in the natural parameters, and
\item running time polynomial in the representation size, $1/\varepsilon$, and $\log(1/\delta)$.
\end{itemize}

This yields four distinct computational notions:

\begin{itemize}
\item efficient realizable proper PAC learning
\item efficient realizable improper PAC learning
\item efficient agnostic proper PAC learning
\item efficient agnostic improper PAC learning
\end{itemize}

\section{Statistical Equivalence Without Runtime Constraints}

At the level of sample complexity, realizable and agnostic PAC learning are equivalent in the distribution-free binary setting. Classically, both are characterized by finite VC dimension; see \citet{blumer1989}. More recently, \citet{hopkins2024} gave a direct black-box reduction explaining realizable-to-agnostic equivalence in a much broader, model-independent way.

This statistical equivalence should not be confused with computational equivalence. The direct reduction of \citet{hopkins2024} is about learnability as a sample-complexity phenomenon. It does \emph{not} by itself settle whether every polynomial-time realizable learner can be converted into a polynomial-time agnostic learner.

\section{The Seven Nontrivial Edges}

Among the $12$ directed implications between the four computational notions above, $5$ are trivial: proper implies improper, agnostic implies realizable, and the obvious compositions of these two facts. The remaining $7$ edges are the ones that need actual theorems, counterexamples, or open-problem markers.

\paragraph{Efficient realizable proper $\Rightarrow$ efficient agnostic proper.}
This implication fails assuming $\mathrm{RP}\neq\mathrm{NP}$. Conjunctions are efficiently properly learnable in the realizable PAC model by the standard elimination algorithm of \citet{valiant1984}. On the other hand, \citet[Theorem~2]{kearns1994} show that properly agnostically learning conjunctions, even when the labels are generated by a polynomial-size DNF, would yield a randomized polynomial-time algorithm for minimum set cover. So realizable proper learning does not in general upgrade to agnostic proper learning.

\paragraph{Efficient realizable proper $\Rightarrow$ efficient agnostic improper.}
No unconditional implication is known, and this edge fails under standard hardness assumptions. Halfspaces are classically efficiently learnable in the realizable proper setting, but \citet{tiegel2023} prove that distribution-free agnostic learning of halfspaces is hard even in the improper setting under worst-case lattice assumptions. Thus realizable proper learnability does not force agnostic improper learnability.

\paragraph{Efficient realizable improper $\Rightarrow$ efficient realizable proper.}
This implication also fails assuming $\mathrm{RP}\neq\mathrm{NP}$. For every fixed $k \ge 2$, \citet[Theorem~3.2]{pitt1988} show that $k$-term DNF is not properly learnable from examples unless $\mathrm{RP}=\mathrm{NP}$. In contrast, \citet[Theorem~3.7]{pitt1988}, building on \citet{valiant1984}, give an improper learner via $k$-CNF, and every fixed-$k$ term DNF has an equivalent polynomial-size $k$-CNF. By \citet{haussler1990}, the examples model used in \citet{pitt1988} is polynomially equivalent to the standard PAC model.

\paragraph{Efficient realizable improper $\Rightarrow$ efficient agnostic proper.}
This edge fails for the same fixed-$k$ DNF witness. If it held, then fixed-$k$ term DNF would be efficiently properly agnostically learnable; restricting that learner to realizable distributions would then give efficient proper realizable learning, contradicting \citet[Theorem~3.2]{pitt1988}. So the fixed-$k$ DNF versus $k$-CNF separation blocks this mixed implication as well.

\paragraph{Efficient realizable improper $\Rightarrow$ efficient agnostic improper.}
This edge is false under the same halfspace hardness used above. Halfspaces are efficiently realizably properly learnable, hence certainly efficiently realizably improperly learnable, but \citet{tiegel2023} rule out even improper agnostic learning under worst-case lattice assumptions. Thus realizable improper learning does not in general upgrade to agnostic improper learning.

\paragraph{Efficient agnostic improper $\Rightarrow$ efficient realizable proper.}
No general theorem or counterexample is currently known in the distribution-free binary polynomial-time setting. \citet[Theorem~3]{kearns1994} show that agnostic learning of one class can already imply realizable PAC learning of a different class via weak approximation; in particular, agnostic learning of conjunctions would imply PAC learning of polynomial-size DNF. At the same time, even for disjunctions, the current best distribution-free agnostic learner is only subexponential-time \citep{diakonikolas2025disjunctions}. We therefore record this edge as open.

\paragraph{Efficient agnostic improper $\Rightarrow$ efficient agnostic proper.}
This implication is also unresolved in general. Existing positive results are class- and distribution-specific rather than black-box properizations: \citet{diakonikolas2021} obtain the first proper agnostic learner for halfspaces under Gaussian marginals, and \citet{lange2023} obtain the first efficient proper agnostic learner for monotone Boolean functions under the uniform distribution while explicitly bypassing a black-box correction barrier. No general theorem or counterexample is currently known for this edge in the distribution-free setting considered here.

\section{Positive Routes To Efficient Agnostic Learning}

\subsection{Efficient Empirical Risk Minimization}

The cleanest route from realizable learning to agnostic learning is \emph{not} a direct learner-to-learner conversion. It is an optimization statement:

\begin{quote}
If empirical $0$-$1$ risk minimization over the benchmark class can be carried out in polynomial time, then agnostic learning follows from uniform convergence together with finite VC dimension.
\end{quote}

This is the right explanation for many classes where realizable and agnostic efficient learning coincide. The key extra ingredient is not realizability itself, but polynomial-time optimization of empirical error.

\subsection{Restricted Distribution Families}

The most common positive results for difficult classes appear once the marginal distribution is restricted. For halfspaces, agnostic learning becomes significantly more tractable under Gaussian, spherical, log-concave, or margin assumptions; see for example \citet{kalai2008agnostically, daniely2015, diakonikolas2021}.

This is exactly why the document should keep ``distribution-free'' and ``restricted-distribution'' cells separate rather than collapsing them into a single agnostic notion.

\section{Barriers And Negative Evidence}

\subsection{The Classical Barrier}

\citet{kearns1994} is still the right starting point for the computational story. Their paper makes clear that agnostic learning is not just a cosmetic weakening of realizability: it is closely tied to harder optimization and noise-tolerance tasks. In particular, the paper already points toward malicious-noise-type difficulties and shows that naive extensions of PAC algorithms need not remain efficient in the agnostic setting.

\subsection{Halfspaces As The Running Example}

Halfspaces illustrate the overall picture well.

\begin{itemize}
\item In the realizable case, halfspaces are efficiently learnable by standard linear-algebraic or linear-programming methods.
\item In the agnostic case, positive results exist under additional structure, such as the uniform sphere or Gaussian marginals; see \citet{daniely2015, diakonikolas2021}.
\item Under stronger hardness assumptions, even improper agnostic learning of halfspaces is computationally hard in the distribution-independent setting; see \citet{tiegel2023}.
\end{itemize}

So for this important class, efficient realizable learnability does \emph{not} by itself deliver an unconditional distribution-free efficient agnostic learner.

\subsection{Properness Is A Real Computational Constraint}

The gap between proper and improper learning is not cosmetic. There are settings where the best known agnostic learner is inherently improper, and later work is needed to recover proper guarantees under extra assumptions. The Gaussian halfspace result of \citet{diakonikolas2021} is a useful anchor point: it is interesting precisely because proper agnostic learning there was missing even though improper agnostic learning was already understood.

\section{Current Verdict For The Atlas}

The first settled conclusion to record is:

\begin{quote}
There is no general theorem currently recorded in this repo saying that every polynomial-time realizable PAC learner can be black-box transformed into a polynomial-time agnostic PAC learner in the distribution-free binary setting.
\end{quote}

What \emph{is} safe to record now is the following:

\begin{itemize}
\item realizable and agnostic PAC learning are statistically equivalent
\item among the $7$ nontrivial computational edges, $5$ are known to fail under standard assumptions
\item the known failures are realizable proper $\Rightarrow$ agnostic proper, realizable proper $\Rightarrow$ agnostic improper, realizable improper $\Rightarrow$ realizable proper, realizable improper $\Rightarrow$ agnostic proper, and realizable improper $\Rightarrow$ agnostic improper
\item the remaining two edges, agnostic improper $\Rightarrow$ realizable proper and agnostic improper $\Rightarrow$ agnostic proper, are currently unresolved in this distribution-free binary polynomial-time regime
\item positive realizable-to-agnostic upgrades usually require extra optimization structure, especially efficient empirical risk minimization or a class-specific argument
\end{itemize}

\section{Multiclass Classification}\label{sec:multiclass}

The binary discussion above does not transfer verbatim to multiclass prediction. Properness already matters at the level of pure sample complexity, and once the label space is unbounded the right combinatorial parameter for learnability is no longer the Natarajan dimension alone. We therefore separate the multiclass story into a sample-complexity subsection and a computational subsection.

\subsection{Sample Complexity}

For bounded label spaces, classical multiclass bounds place learnability between the Natarajan and graph dimensions, and \citet{daniely2015multiclass} show that different ERM learners can differ by a factor of $\Theta(\log |Y|)$ in sample complexity. In full generality, \citet{brukhim2022} prove that realizable PAC learnability is characterized by finite DS dimension, while \citet{cohen2025multiclass} show that the current best general agnostic bounds depend on both the DS dimension $d_{DS}$ and the Natarajan dimension $d_N$, scaling up to logarithmic factors like $\widetilde{O}(d_{DS}^{3/2}/\varepsilon + d_N/\varepsilon^2)$.

The same four nodes are relevant here: sample-efficient realizable proper, sample-efficient realizable improper, sample-efficient agnostic proper, and sample-efficient agnostic improper. As before, there are $12$ directed edges, $5$ of which are trivial. The remaining $7$ behave differently from the binary case.

\paragraph{Realizable proper $\Rightarrow$ agnostic proper.}
This implication is true. The reduction of \citet{hopkins2024} is model-independent and applies equally to proper multiclass learning problems, so any sample-efficient proper learner in the realizable setting yields a sample-efficient proper learner in the agnostic setting as well.

\paragraph{Realizable proper $\Rightarrow$ agnostic improper.}
This implication is also true. It follows from the previous paragraph and then forgetting the properness requirement. So realizability still suffices for agnostic learnability at the level of sample complexity.

\paragraph{Realizable improper $\Rightarrow$ realizable proper.}
This implication fails. \citet[Theorem~1]{daniely2014optimal} construct the first Cantor classes $H_d$ and $H_\infty$. The finite classes $H_d$ already show a large proper-versus-improper sample gap. The limit class $H_\infty$ is PAC learnable but not learnable by any proper algorithm. Thus properness is a genuine statistical constraint in multiclass learning.

\paragraph{Realizable improper $\Rightarrow$ agnostic proper.}
This implication fails for the same $H_\infty$ witness. If improper realizable learnability implied proper agnostic learnability, then restricting the agnostic learner to realizable distributions would give a proper realizable learner, contradicting \citet[Theorem~1]{daniely2014optimal}.

\paragraph{Realizable improper $\Rightarrow$ agnostic improper.}
This implication is true. It is the multiclass improper instance of the realizable-to-agnostic reduction of \citet{hopkins2024}. In particular, once a multiclass class is statistically learnable in the realizable improper sense, it is statistically learnable in the agnostic improper sense as well.

\paragraph{Agnostic improper $\Rightarrow$ realizable proper.}
This implication fails. Since $H_\infty$ is PAC learnable, \citet{brukhim2022} imply that it has finite DS dimension and is therefore agnostically learnable in the improper model. But \citet[Theorem~1]{daniely2014optimal} show that no proper learner can learn $H_\infty$ even in the realizable case.

\paragraph{Agnostic improper $\Rightarrow$ agnostic proper.}
This implication fails as well. The same $H_\infty$ witness suffices: improper agnostic learning is available by \citet{brukhim2022}, whereas proper agnostic learning would immediately imply proper realizable learning, which \citet[Theorem~1]{daniely2014optimal} rule out.

\subsection{Computational Complexity}

At the polynomial-time level, the coarse implication graph does not improve in the multiclass setting, because binary classification is a special case of multiclass classification. Every binary counterexample from the previous section therefore lifts verbatim to multiclass by viewing the same class as taking values in a label set with at least two labels. Consequently, the status of the seven nontrivial edges is the same as in the binary case.

\paragraph{Efficient realizable proper $\Rightarrow$ efficient agnostic proper.}
False under standard assumptions already for the binary subclass of conjunctions: realizable proper learning is efficient, but proper agnostic learning is hard unless $\mathrm{RP}=\mathrm{NP}$ \citep{valiant1984,kearns1994}.

\paragraph{Efficient realizable proper $\Rightarrow$ efficient agnostic improper.}
False under standard hardness assumptions already for the binary subclass of halfspaces: realizable proper learning is efficient, but even improper agnostic learning is hard under worst-case lattice assumptions \citep{tiegel2023}.

\paragraph{Efficient realizable improper $\Rightarrow$ efficient realizable proper.}
False already for the binary subclass of fixed-$k$ term DNF. \citet{pitt1988} give an improper realizable learner via $k$-CNF while also showing that proper learning is hard unless $\mathrm{RP}=\mathrm{NP}$; see also \citet{haussler1990}.

\paragraph{Efficient realizable improper $\Rightarrow$ efficient agnostic proper.}
False for the same fixed-$k$ DNF witness. If this implication held in multiclass, it would in particular hold for the binary subclass and contradict the proper-learning hardness of \citet{pitt1988}.

\paragraph{Efficient realizable improper $\Rightarrow$ efficient agnostic improper.}
False for the same halfspace witness as above. The improper agnostic hardness from \citet{tiegel2023} already rules out this implication in the multiclass setting.

\paragraph{Efficient agnostic improper $\Rightarrow$ efficient realizable proper.}
Still open in general. Any theorem or counterexample here would immediately settle the corresponding binary edge, and no such general result is currently known. The literature instead gives class-specific properization results and partial algorithms rather than a universal implication \citep{diakonikolas2021,lange2023,diakonikolas2025disjunctions}.

\paragraph{Efficient agnostic improper $\Rightarrow$ efficient agnostic proper.}
This edge is likewise open in general for the same reason: a multiclass resolution would specialize to the unresolved binary problem. Known progress remains class-specific rather than a black-box properization theorem \citep{diakonikolas2021,lange2023}.

There are nevertheless positive multiclass-specific efficient results for concrete classes. For generalized linear multiclass predictors, \citet{daniely2014optimal} give efficient improper learners with optimal sample complexity in the realizable case and, via compression-based arguments, in the agnostic case as well. These class-specific algorithms do not change the global implication graph above, but they show that the obstructions are not universal across all multiclass classes.

\section{Regression With Continuum Labels}

To keep the regression discussion precise, this section fixes the standard bounded regression model: labels lie in $[0,1]$, hypotheses are functions $f : \mathcal{X} \to [0,1]$, and performance is measured by squared loss $(f(x)-y)^2$. This is the classical real-valued PAC or agnostic regression setting of \citet{bartlett1996, alon1997}. If later we want to track absolute loss, unbounded labels, or other loss functions, those should be split into separate cells.

\subsection{Sample Complexity}

For bounded real-valued regression, the statistical story is much cleaner than in multiclass classification. \citet{bartlett1996} and \citet{alon1997} show that learnability is governed by the scale-sensitive fat-shattering or $P_\gamma$ dimensions, and the learner in these classical formulations is already \emph{proper}: it outputs a function from the benchmark class itself. Moreover, \citet{hopkins2024} explicitly include regression in their realizable-to-agnostic reduction, so realizable and agnostic learnability remain equivalent at the sample-complexity level.

As a result, in this standard bounded-loss regression model all seven nontrivial statistical edges are true.

\paragraph{Realizable proper $\Rightarrow$ agnostic proper.}
True by the realizable-to-agnostic reduction of \citet{hopkins2024}, which applies to regression and more general loss-based settings.

\paragraph{Realizable proper $\Rightarrow$ agnostic improper.}
True by the same reduction and then forgetting properness.

\paragraph{Realizable improper $\Rightarrow$ realizable proper.}
True. If a regression class is learnable at all, then the classical characterizations of \citet{bartlett1996, alon1997} provide a proper learner for that same class, so improper learning gives no extra statistical power here.

\paragraph{Realizable improper $\Rightarrow$ agnostic proper.}
True by combining the previous paragraph with \citet{hopkins2024}.

\paragraph{Realizable improper $\Rightarrow$ agnostic improper.}
True, again by \citet{hopkins2024}.

\paragraph{Agnostic improper $\Rightarrow$ realizable proper.}
True. Agnostic learnability implies learnability of the class, and in bounded regression learnability already entails a proper learner by \citet{bartlett1996, alon1997}.

\paragraph{Agnostic improper $\Rightarrow$ agnostic proper.}
True for the same reason: once the class is statistically learnable, the classical proper learning results recover a proper agnostic learner.

\subsection{Computational Complexity}

At the polynomial-time level, regression again loses the clean statistical collapse. For natural convex classes such as linear predictors under squared loss, empirical risk minimization is convex and efficient realizable and agnostic proper learning coincide. But for the general implication graph, the binary computational barriers from earlier in the chapter transfer directly.

The transfer is elementary. Given a binary class $\mathcal{C} \subseteq \{0,1\}^{\mathcal{X}}$, view it as a regression class $\mathcal{F} = \mathcal{C} \subseteq [0,1]^{\mathcal{X}}$ and use squared loss on labels in $\{0,1\}$. For proper learners, squared loss is exactly $0$-$1$ loss on this subclass. For improper learners, thresholding a predictor at $1/2$ converts a squared-loss guarantee into a classification guarantee up to a constant factor. Therefore every binary computational separation and every binary open edge carries over to regression.

So the seven nontrivial computational edges have the same status as in the binary section:

\begin{itemize}
\item the five edges realizable proper $\Rightarrow$ agnostic proper, realizable proper $\Rightarrow$ agnostic improper, realizable improper $\Rightarrow$ realizable proper, realizable improper $\Rightarrow$ agnostic proper, and realizable improper $\Rightarrow$ agnostic improper all fail under the same standard assumptions and witnesses as before \citep{valiant1984,kearns1994,pitt1988,haussler1990,tiegel2023}
\item the two edges agnostic improper $\Rightarrow$ realizable proper and agnostic improper $\Rightarrow$ agnostic proper remain open in the same generality for the same reason: any resolution for regression would immediately resolve the corresponding binary problem
\end{itemize}

Thus regression provides the opposite contrast to multiclass classification. Statistically, the bounded real-valued theory is cleaner than multiclass and behaves more like the classical binary fundamental theorem. Computationally, however, the broad improper-versus-proper and realizable-versus-agnostic implication graph is no easier in full generality.

\section{Open Cells To Track Next}

This chapter should expand next along the following axes:

\begin{itemize}
\item find the cleanest theorem statement linking efficient proper realizable learning to efficient proper agnostic learning under explicit ERM assumptions
\item add a dedicated subsection on malicious noise as an intermediate notion between realizable and agnostic learning
\item separate unconditional separations from cryptographic or average-case hardness evidence
\item add at least one non-halfspace example where realizable and agnostic efficient learning provably coincide for structural reasons
\item decide whether to treat efficient weak agnostic learning and boosting as a separate bridge chapter
\end{itemize}

\section{Appendix: Proof Sketches For Chapter 2}

This appendix records proof ideas for the main claims in this chapter. The goal is not to reproduce the full arguments from the cited papers, but to isolate the key reductions, combinatorial parameters, and estimates so that a strong undergraduate reader can reconstruct the proofs with the original sources in hand.

\subsection{Binary Classification}

\paragraph{Why efficient ERM gives efficient agnostic learning.}
This is the cleanest positive proof in the chapter. Fix a class $\mathcal{C}$ of finite VC dimension and assume we can compute an empirical risk minimizer over $\mathcal{C}$ in polynomial time. Uniform convergence says that for a sample of size
\[
n = O\!\left(\frac{\mathrm{VC}(\mathcal{C}) + \log(1/\delta)}{\varepsilon^2}\right),
\]
every hypothesis has empirical error within about $\varepsilon$ of its true error. So if $\hat h$ minimizes empirical error, then
\[
\err(\hat h) \le \widehat{\err}(\hat h) + \varepsilon \le \widehat{\err}(h^\star) + \varepsilon \le \err(h^\star) + 2\varepsilon,
\]
where $h^\star$ is the best hypothesis in the class. The whole proof is just this three-line inequality plus the VC uniform-convergence theorem. The only computational ingredient is the ability to find $\hat h$ efficiently.

\paragraph{Why conjunctions separate realizable proper from agnostic proper.}
First fix the learning problem. Let $\mathcal{T}_n$ be the class of conjunctions of literals over variables $x_1,\dots,x_n$. In the realizable proper PAC problem, the labels are generated by an unknown target $c \in \mathcal{T}_n$, and we must output some $h \in \mathcal{T}_n$ with small error. In the agnostic proper problem, the examples come from an arbitrary distribution on $\{0,1\}^n \times \{0,1\}$, and we must output $h \in \mathcal{T}_n$ with error at most
\[
\opt(\mathcal{T}_n) + \varepsilon.
\]

For the realizable side, the clean proof really is ``finite VC dimension plus efficient ERM.'' First, $\mathcal{T}_n$ has finite VC dimension. In fact one can show
\[
\mathrm{VC}(\mathcal{T}_n)=n,
\]
but finiteness is enough here. The easy upper bound is by counting: for each variable $x_i$, a conjunction can either include $x_i$, include $\neg x_i$, or omit $x_i$, so
\[
|\mathcal{T}_n| \le 3^n.
\]
Any class that shatters $d$ points must have at least $2^d$ distinct labelings on those points, hence
\[
2^d \le |\mathcal{T}_n| \le 3^n,
\]
so $d \le n \log_2 3$. Therefore $\mathcal{T}_n$ has finite VC dimension. If you want the matching lower bound, take the $n$ points
\[
x^{(i)} = (1,\dots,1,0,1,\dots,1),
\]
with the zero in coordinate $i$. Any subset $S \subseteq [n]$ is realized by the monotone conjunction $\bigwedge_{i \notin S} x_i$, so these $n$ points are shattered.

Once finite VC dimension is known, the realizable PAC theorem says that any consistent ERM learner for $\mathcal{T}_n$ has polynomial sample complexity. So the only remaining job is computational: show that ERM, or at least finding a consistent conjunction, is easy. The standard learner is the elimination algorithm, and yes, it can be viewed as a very concrete ERM or consistency algorithm. Start with the conjunction of all $2n$ literals
\[
x_1,\neg x_1,\dots,x_n,\neg x_n,
\]
and whenever a positive example arrives, delete every literal it falsifies. The final hypothesis is still a conjunction of literals, and the whole procedure runs in time $O(mn)$ on $m$ examples. Every target literal survives, because every positive example satisfies the target conjunction. Hence every negative example is also rejected by the learned conjunction, since any negative example must violate at least one target literal that is still present. So in the realizable setting the algorithm outputs a hypothesis with zero empirical error, which is exactly an ERM solution. Combining this with the finite-VC argument gives efficient proper realizable PAC learning of conjunctions.

For the hardness proof, \citet{kearns1994} reduce minimum set cover to proper agnostic learning of conjunctions. Start from a set cover instance with objects $o_1,\dots,o_t$ and sets $S_1,\dots,S_n$. The variables $x_1,\dots,x_n$ correspond to the sets. The target labeling function is the polynomial-size DNF
\[
f = T_1 \vee \cdots \vee T_n,
\qquad
T_j = \bigwedge_{k \ne j} x_k.
\]
So $f(x)=1$ exactly when the assignment has at most one zero coordinate.

Now build three kinds of labeled points. For each object $o_i$, define a point $a_i \in \{0,1\}^n$ by setting
\[
(a_i)_j =
\begin{cases}
0 & \text{if } o_i \in S_j,\\
1 & \text{otherwise.}
\end{cases}
\]
Because every object is assumed to belong to at least two sets, each $a_i$ has at least two zero coordinates, so $f(a_i)=0$. These are the negative examples. For each set $S_j$, define $b_j$ to be the point with a single zero in coordinate $j$ and ones elsewhere; then $f(b_j)=1$. Finally let $c=(1,\dots,1)$, which is also positive.

The key dictionary is this: if we use a \emph{monotone} conjunction
\[
h_B(x)=\bigwedge_{j \in B} x_j,
\]
then choosing the variable $x_j$ means choosing the set $S_j$. Indeed, $h_B(a_i)=0$ if and only if some chosen coordinate $j \in B$ satisfies $(a_i)_j=0$, which happens exactly when $o_i \in S_j$. So $h_B$ rejects every negative point $a_i$ exactly when the chosen sets $B$ form a set cover.

The distribution is then chosen to force three behaviors. First, the point $c$ gets probability mass $1/2$. This makes any negative literal $\neg x_j$ impossible in a near-optimal hypothesis, because such a literal misclassifies $c$ all by itself and already incurs error about $1/2$, while there is an obvious monotone conjunction with much smaller error. So every near-optimal conjunction must be monotone. Second, each negative point $a_i$ gets fairly large mass, on the order of $1/(t+1)$. Therefore any near-optimal conjunction must classify every $a_i$ correctly, so the selected variables must form a set cover. Third, each positive point $b_j$ gets tiny mass
\[
\frac{1}{4n(t+1)}.
\]
Including the literal $x_j$ causes exactly one extra positive mistake, namely on $b_j$. Therefore, once we know the learner has to output a monotone conjunction that covers all the $a_i$, its total error is exactly proportional to the number of selected sets:
\[
\err(h_B) = \frac{|B|}{4n(t+1)}.
\]

That is the whole reduction. Near-optimal error means: cover all objects, and among all covers choose one with minimum cardinality. By setting
\[
\varepsilon = \frac{1}{8n(t+1)},
\]
Kearns, Schapire, and Sellie ensure that an additive-$\varepsilon$ agnostic learner must recover an \emph{exact} minimum set cover, since changing the cover size by even one changes the error by
\[
\frac{1}{4n(t+1)} > \varepsilon.
\]
So an efficient proper agnostic learner for conjunctions would yield an efficient algorithm for minimum set cover.

\paragraph{Why fixed-$k$ term DNF separates realizable improper from realizable proper.}
The positive side is the easier part. If
\[
f = T_1 \vee T_2 \vee \cdots \vee T_k
\]
is a DNF with fixed $k$, then distributing $\vee$ over $\wedge$ converts it into a logically equivalent $k$-CNF whose size is polynomial in the original description because $k$ is constant. Once you write down the distribution identity for $k=2$, the general case is clear. So Valiant's learner for fixed-$k$ CNF gives an improper learner for fixed-$k$ term DNF.

The negative side in \citet{pitt1988} is a reduction from an $\mathrm{NP}$-hard problem to proper learning from examples. The conceptual point is that requiring the learner to return a DNF of the same restricted syntactic form forces it to solve a hidden combinatorial search problem, while the improper learner escapes that search by moving to the equivalent CNF representation. That is exactly the kind of phenomenon this atlas wants to isolate: representational freedom can matter even before agnosticity enters.

\paragraph{Why halfspaces witness realizable-easy but agnostic-hard.}
For realizable learning, the proof idea is geometric. If a labeled sample is separated by some hyperplane, then one can either run the perceptron algorithm or solve a feasibility linear program to produce a separating halfspace. Both are polynomial-time procedures in the usual finite-dimensional bit-model.

The hardness result of \citet{tiegel2023} runs in the opposite direction. It builds distributions over labeled examples from lattice instances so that, in one case, some halfspace has very small error, while in the other case every hypothesis has noticeably larger error. An agnostic learner that achieved inverse-polynomial excess error would distinguish these two cases, which would in turn solve the lattice problem. For a reader reconstructing the proof, the key task is to track exactly how the lattice promise becomes an error gap for halfspaces.

\paragraph{Why the two binary open edges remain open.}
The unresolved edges
\[
\text{agnostic improper} \Rightarrow \text{realizable proper}
\quad\text{and}\quad
\text{agnostic improper} \Rightarrow \text{agnostic proper}
\]
would require either a generic properization theorem or a concrete witness class with an efficient improper agnostic learner but no efficient proper learner. The natural candidate witness is conjunctions or disjunctions: proper agnostic learning is hard, but efficient improper agnostic learning is itself still unresolved. That is why the literature currently gives pressure in both directions without closing either edge.

\subsection{Multiclass Classification}

\paragraph{Why realizable-to-agnostic remains true statistically.}
This is the easy part of the multiclass section. Once the learning problem is fixed, \citet{hopkins2024} give a black-box conversion from realizable to agnostic learning. Nothing in the reduction uses the fact that labels are binary; it only needs access to a learner for the underlying hypothesis class. So the realizable proper and realizable improper nodes both inherit their agnostic counterparts at the sample-complexity level.

\paragraph{Why the Cantor classes separate proper from improper.}
The first Cantor class is indexed by subsets $A$ of a finite set. Each hypothesis $h_A$ gives a special label on points in $A$ and a default label on points outside $A$. The improper learner only needs to predict well relative to this family, not to recover the exact set $A$, and that freedom lets one compress the relevant information in the sample very efficiently. A proper learner, by contrast, must output some actual $h_A$, so it has to determine much more of the hidden subset structure. That is why the sample complexity jumps.

For the infinite Cantor class $H_\infty$, the same idea becomes absolute rather than quantitative. Improper learning is still possible because the class remains statistically learnable, but proper learning fails altogether because any proper learner can be tricked on an unseen point whose membership in the hidden subset was never revealed. The proof is worth reconstructing directly from the definition of the class: one simply writes down a distribution concentrated on finitely many unseen possibilities and shows that every proper output makes a forced mistake on at least one of them.

\paragraph{Why the multiclass computational edges match the binary ones.}
This proof is just an embedding. Start with a binary class $\mathcal{C}$ and view it as a multiclass class whose label set merely happens to have two labels. Any multiclass algorithm for that class would immediately be a binary algorithm as well. Therefore every binary counterexample, hardness result, and open edge automatically survives in the multiclass setting. The only thing to check is that the definitions of properness, agnosticity, and efficient learning are literally the same after the relabeling.

\subsection{Regression With Continuum Labels}

\paragraph{Why all seven statistical edges collapse in bounded regression.}
Here the main theorem is not about classification at all, but about the structure of real-valued function classes. The classical results of \citet{bartlett1996, alon1997} show that bounded regression is learnable exactly when its fat-shattering dimensions are finite at all scales, and the learner they construct is proper. So once a class is statistically learnable in any sense, a proper learner already exists.

That immediately kills the proper-versus-improper distinction at the sample-complexity level. Then \citet{hopkins2024} supply the realizable-to-agnostic conversion. Putting the two facts together gives all seven nontrivial edges. A good exercise is to write these implications as a commutative diagram: one theorem collapses proper versus improper, and the other collapses realizable versus agnostic.

\paragraph{Why binary hardness transfers to regression.}
This reduction is elementary and useful enough to remember. Take a binary class $\mathcal{C} \subseteq \{0,1\}^{\mathcal{X}}$ and regard it as a regression class $\mathcal{F} \subseteq [0,1]^{\mathcal{X}}$ under squared loss. On labels $y \in \{0,1\}$ and proper hypotheses $h \in \{0,1\}^{\mathcal{X}}$, the squared loss
\[
(h(x)-y)^2
\]
is exactly the $0$-$1$ loss. So every proper classification lower bound is already a proper regression lower bound.

For improper learners, let $r(x) \in [0,1]$ be the real-valued prediction and threshold it at $1/2$:
\[
\tilde h(x) = \mathbf{1}[r(x) \ge 1/2].
\]
If $\tilde h(x) \neq y$, then $|r(x)-y| \ge 1/2$, hence $(r(x)-y)^2 \ge 1/4$. Therefore
\[
\Pr[\tilde h(x) \neq y] \le 4 \, \mathbb{E}(r(x)-y)^2.
\]
This one-line inequality is the whole reduction. Any efficient improper regression learner with small excess squared loss would therefore give an efficient improper classifier with comparable error, so the binary computational separations and open problems carry over unchanged.
